\documentclass[11pt,a4paper]{article}

% --- Encoding ---
\usepackage[utf8]{inputenc}
\usepackage[T1]{fontenc}
\usepackage[spanish,es-noshorthands]{babel}
\usepackage{lmodern}

% --- Layout ---
\usepackage[margin=2.5cm]{geometry}
\usepackage{parskip}
\setlength{\parindent}{0pt}
\setlength{\parskip}{0.5em plus 0.2em minus 0.1em}

% --- Tables and graphics ---
\usepackage{booktabs}
\usepackage{array}
\usepackage{enumitem}
\usepackage{xcolor}
\usepackage{graphicx}
\usepackage{tikz}
\usetikzlibrary{positioning,shapes.geometric,arrows.meta,fit,backgrounds,calc,decorations.pathreplacing}
\usepackage{caption}
\usepackage{subcaption}
\usepackage{amsmath,amssymb}

% --- Soporte para caracteres Unicode usados en fórmulas en texto plano ---
\usepackage{newunicodechar}
\newunicodechar{→}{\ensuremath{\rightarrow}}
\newunicodechar{≥}{\ensuremath{\geq}}
\newunicodechar{≤}{\ensuremath{\leq}}
\newunicodechar{×}{\ensuremath{\times}}
\newunicodechar{÷}{\ensuremath{\div}}
\newunicodechar{≈}{\ensuremath{\approx}}
\newunicodechar{⌈}{\ensuremath{\lceil}}
\newunicodechar{⌉}{\ensuremath{\rceil}}
\newunicodechar{✓}{\ensuremath{\checkmark}}

% --- Code listings ---
\usepackage{listings}
\definecolor{codegray}{rgb}{0.96,0.96,0.96}
\definecolor{codeborder}{rgb}{0.85,0.85,0.85}
\definecolor{codecomment}{rgb}{0.30,0.55,0.30}
\definecolor{codekeyword}{rgb}{0.10,0.30,0.70}
\definecolor{codestring}{rgb}{0.65,0.15,0.10}

\lstdefinestyle{cstyle}{
    language=C,
    backgroundcolor=\color{codegray},
    basicstyle=\ttfamily\footnotesize,
    breaklines=true,
    frame=single,
    rulecolor=\color{codeborder},
    showstringspaces=false,
    keywordstyle=\color{codekeyword}\bfseries,
    commentstyle=\color{codecomment}\itshape,
    stringstyle=\color{codestring},
    aboveskip=1em,
    belowskip=1em,
}

\lstdefinestyle{shellstyle}{
    language=bash,
    backgroundcolor=\color{codegray},
    basicstyle=\ttfamily\small,
    breaklines=true,
    frame=single,
    rulecolor=\color{codeborder},
    showstringspaces=false,
    keywordstyle=\color{codekeyword}\bfseries,
    commentstyle=\color{codecomment}\itshape,
    stringstyle=\color{codestring},
    aboveskip=1em,
    belowskip=1em,
    columns=fullflexible,
    upquote=true,
}

\lstdefinestyle{textstyle}{
    backgroundcolor=\color{codegray},
    basicstyle=\ttfamily\footnotesize,
    breaklines=true,
    frame=single,
    rulecolor=\color{codeborder},
    showstringspaces=false,
    aboveskip=1em,
    belowskip=1em,
    columns=fullflexible,
}

\lstset{style=textstyle}

% --- Definition / Idea / Important boxes ---
\usepackage[most]{tcolorbox}
\newtcolorbox{definicion}[1][]{
    colback=blue!5!white,
    colframe=blue!50!black,
    fonttitle=\bfseries,
    title=Definición,
    sharp corners=southwest,
    rounded corners=northwest,
    arc=2pt,
    boxrule=0.6pt,
    left=8pt,right=8pt,top=4pt,bottom=4pt,
    #1
}
\newtcolorbox{idea}[1][]{
    colback=orange!5!white,
    colframe=orange!60!black,
    fonttitle=\bfseries,
    title=Idea intuitiva,
    sharp corners=southwest,
    rounded corners=northwest,
    arc=2pt,
    boxrule=0.6pt,
    left=8pt,right=8pt,top=4pt,bottom=4pt,
    #1
}
\newtcolorbox{importante}[1][]{
    colback=red!5!white,
    colframe=red!50!black,
    fonttitle=\bfseries,
    title=Importante,
    sharp corners=southwest,
    rounded corners=northwest,
    arc=2pt,
    boxrule=0.6pt,
    left=8pt,right=8pt,top=4pt,bottom=4pt,
    #1
}

% --- Hyperlinks ---
\usepackage[colorlinks=true,
            linkcolor=blue!60!black,
            urlcolor=blue!50!black,
            citecolor=blue!60!black]{hyperref}

% --- Color palette ---
\definecolor{io1}{HTML}{4C72B0}
\definecolor{io2}{HTML}{DD8452}
\definecolor{io3}{HTML}{55A868}
\definecolor{io4}{HTML}{8172B2}
\definecolor{ioused}{HTML}{C44E52}
\definecolor{iofree}{HTML}{8C8C8C}
\definecolor{ioshared}{HTML}{4FC3C7}
\definecolor{iofile}{HTML}{F4A261}

% --- Title block ---
\title{
    \vspace{-2.5em}
    \textbf{Práctico 6} \\[0.3em]
    \large Entrada/salida, filesystem y usuarios
}
\author{
    Tomás Rodeghiero \\
    \small Sistemas Operativos --- Universidad Nacional de Río Cuarto
}
\date{2026}

\begin{document}

\maketitle

\begin{abstract}
\noindent
Este documento presenta la resolución integral del Práctico 6 de la
materia Sistemas Operativos (UNRC, 2026), centrado en
\emph{entrada/salida}, \emph{sistemas de archivos} y
\emph{usuarios/control de acceso}. Se resuelven los \textbf{diez
ejercicios} de la consigna oficial, divididos en dos bloques. El
\textbf{bloque A} recorre los conceptos de E/S y filesystem: dispositivos
de bloque y carácter en \texttt{/dev}, uso de \texttt{dd} y de los
dispositivos especiales \texttt{/dev/zero} y \texttt{/dev/urandom}, el
flujo de compilación y carga de un módulo del kernel (RTC), la
representación de un directorio en FAT vs.\ inodos, y la API
\texttt{stat}/\texttt{readdir} para navegar el filesystem desde C. El
\textbf{bloque B} aborda usuarios y permisos: listado de usuarios y
grupos en Linux y macOS, semántica de permisos UGO/octal (con énfasis
en el bit de ejecución sobre directorios), los UIDs real y efectivo, y
el funcionamiento del bit \texttt{SUID} a través de un ejemplo
ejecutable paso a paso. Cada ejercicio se justifica con la teoría del
capítulo \emph{Entrada/salida y sistemas de archivos} de las Notas del
curso (Marcelo Arroyo).
\end{abstract}

\tableofcontents
\bigskip

% =====================================================================
\section{Introducción}
% =====================================================================

El subsistema de \emph{entrada/salida} es la capa del SO que abstrae el
hardware heterogéneo de los dispositivos y lo presenta como una
\textbf{API uniforme} para los procesos. La filosofía clásica de UNIX
---``todo es un archivo''--- llevó esta idea hasta sus últimas
consecuencias: discos, teclados, ratones, audio, red, dispositivos
especiales (\texttt{/dev/zero}, \texttt{/dev/urandom},
\texttt{/dev/null}) y hasta el propio kernel (\texttt{/proc},
\texttt{/sys}) se acceden con las mismas syscalls
(\texttt{open}/\texttt{read}/\texttt{write}/\texttt{close}/\texttt{ioctl}).

\begin{center}
\begin{tikzpicture}[
    every node/.style={font=\small},
    layer/.style={draw,minimum width=10cm,minimum height=0.8cm,align=center}
]
\node[layer,fill=io1!20]                       (api)   at (0, 3.6) {\textbf{API de usuario} (\texttt{read}, \texttt{write}, \texttt{open}, \texttt{ioctl}, \texttt{mmap})};
\node[layer,fill=io2!20,below=0.05cm of api]   (kio)   {\textbf{Kernel I/O subsystem} (buffering, caching, spooling, scheduling)};
\node[layer,fill=io3!20,below=0.05cm of kio]   (drv)   {\textbf{Drivers} (uno por familia de dispositivo)};
\node[layer,fill=io4!20,below=0.05cm of drv]   (hw)    {\textbf{Hardware} (controladores, registros, DMA, IRQs)};
\end{tikzpicture}
\end{center}

\textbf{Modos de transferencia.} El SO puede mover datos entre la
memoria y un dispositivo de tres formas:

\begin{itemize}[leftmargin=*,nosep]
    \item \textbf{Programmed I/O (PIO)}: la CPU mueve cada byte
        leyendo/escribiendo registros del controlador. Es simple pero
        gasta toda la CPU del proceso.
    \item \textbf{Interrupt-driven I/O}: la CPU le pide al controlador
        que haga la operación y se va a hacer otra cosa. Cuando el
        controlador termina, dispara una \emph{IRQ}.
    \item \textbf{DMA} (\emph{Direct Memory Access}): un controlador
        independiente mueve datos entre el dispositivo y la RAM sin
        intervención de la CPU. La CPU solo participa al inicio
        (configurar la transferencia) y al final (recibir la IRQ de
        terminación).
\end{itemize}

\begin{importante}
DMA \textbf{no} libera completamente a la CPU. La CPU sigue siendo
quien arranca la transferencia (escribe en los registros del DMA
controller) y quien atiende la IRQ cuando termina. Lo que se ahorra es
el ciclo byte-a-byte intermedio, que es donde está el grueso del
trabajo en E/S de bloque.
\end{importante}

\textbf{Tipos de dispositivos.}

\begin{itemize}[leftmargin=*,nosep]
    \item \textbf{Dispositivos de bloque}: leen y escriben en
        \emph{bloques} de tamaño fijo (típicamente $512$\,B o
        $4$\,KiB). Permiten acceso aleatorio (\texttt{lseek}).
        Discos, SSDs, particiones. Aparecen con la letra \texttt{b}
        en \texttt{ls -l}.
    \item \textbf{Dispositivos de carácter}: leen y escriben byte a
        byte, normalmente sin posibilidad de \texttt{seek}. Teclados,
        seriales, terminales, \texttt{/dev/null}, \texttt{/dev/zero},
        \texttt{/dev/urandom}. Aparecen con \texttt{c}.
    \item \textbf{Dispositivos de red}: no aparecen en \texttt{/dev};
        se acceden con la API de sockets, que es un modelo distinto
        (\texttt{socket}, \texttt{bind}, \texttt{connect}, \ldots).
\end{itemize}

% =====================================================================
\part*{Bloque A --- Entrada/salida y filesystem}
\addcontentsline{toc}{section}{Bloque A --- Entrada/salida y filesystem}
% =====================================================================

% =====================================================================
\section{Ejercicio 1 --- Regex para dispositivos de bloque en \texttt{/dev}}
% =====================================================================

\textbf{Enunciado.} Dar la expresión regular en el comando
\verb|ls -l /dev | grep "regexp here"| para que liste \emph{sólo} los
dispositivos de bloques.

\subsection*{Análisis}

\texttt{ls -l} imprime una línea por entrada con el formato

\begin{lstlisting}[style=textstyle]
<tipo><permisos> <links> <user> <group> <maj>,<min> <fecha> <nombre>
\end{lstlisting}

El \textbf{primer carácter} de cada línea indica el tipo del archivo:

\begin{center}
\small
\begin{tabular}{@{}clp{8.5cm}@{}}
\toprule
Letra & Tipo & Descripción \\
\midrule
\texttt{-} & regular         & archivo común (texto, binario, etc.) \\
\texttt{d} & directorio       & contiene entradas (nombre $\to$ inodo) \\
\texttt{l} & symlink          & contiene una ruta hacia otro archivo \\
\texttt{b} & block device     & dispositivo accedido en bloques (discos) \\
\texttt{c} & character device & dispositivo accedido byte a byte (tty, urandom) \\
\texttt{p} & FIFO / pipe      & cola con nombre creada con \texttt{mkfifo} \\
\texttt{s} & socket           & socket de dominio UNIX \\
\bottomrule
\end{tabular}
\end{center}

Por lo tanto, listar los \emph{dispositivos de bloque} se reduce a
filtrar líneas cuyo primer carácter sea \texttt{b}. La expresión
correcta es \verb|^b|: el \verb|^| ancla el match al \textbf{principio
de la línea}.

\subsection*{Solución}

\begin{lstlisting}[style=shellstyle]
ls -l /dev | grep "^b"
\end{lstlisting}

\textbf{Salida típica en Linux.}

\begin{lstlisting}[style=textstyle]
brw-rw----  1 root disk    8,   0 ... sda
brw-rw----  1 root disk    8,   1 ... sda1
brw-rw----  1 root disk    8,   2 ... sda2
brw-rw----  1 root cdrom  11,   0 ... sr0
brw-rw----  1 root disk    7,   0 ... loop0
\end{lstlisting}

\textbf{Salida típica en macOS.}

\begin{lstlisting}[style=textstyle]
brw-r-----  1 root operator   1,   0 ... disk0
brw-r-----  1 root operator   1,   1 ... disk0s1
brw-r-----  1 root operator   1,   2 ... disk0s2
\end{lstlisting}

\begin{idea}
La línea \texttt{total N} que \texttt{ls -l} imprime al principio del
listado empieza con \texttt{t}, así que \verb|^b| ya la descarta. Si
se omitiera el ancla \verb|^| y se usara solo \verb|"b"|, también
matchearían líneas cuyo nombre contiene una \texttt{b} (\texttt{sda1}
no, pero un dispositivo llamado \texttt{cuab0} por ejemplo sí), lo que
introduciría falsos positivos. El \verb|^| evita exactamente eso.
\end{idea}

\textbf{Conclusión.} La expresión regular pedida es \texttt{"\textasciicircum b"}.

% =====================================================================
\section{Ejercicio 2 --- Archivo de 64 MB lleno de ceros con \texttt{dd}}
% =====================================================================

\textbf{Enunciado.} Crear un archivo de $64$\,MB inicializado en cero,
usando \texttt{dd} y leyendo del dispositivo \texttt{/dev/zero}.

\subsection*{Análisis}

\texttt{dd} (\emph{disk dump / data definition}) copia bytes de
entrada en bytes de salida con tamaño y cantidad de bloques
configurables:

\begin{lstlisting}[style=shellstyle]
dd if=<entrada> of=<salida> bs=<tamano_bloque> count=<nro_bloques>
\end{lstlisting}

El total escrito es exactamente \texttt{bs × count}. \texttt{/dev/zero}
es un dispositivo de \emph{carácter} especial que provee bytes nulos
($\texttt{0x00}$) infinitos: cada \texttt{read} sobre él retorna
buffers llenos de ceros, sin tocar el disco.

Para escribir $64$\,MB tenemos varias combinaciones equivalentes:

\quad $64 \text{\,MB} = 64 × 2^{20}$ bytes = $67\,108\,864$ bytes

\begin{center}
\small
\begin{tabular}{@{}lll@{}}
\toprule
\texttt{bs} & \texttt{count} & total \\
\midrule
\texttt{1M}    & $64$         & $64$\,MB \\
\texttt{4K}    & $16\,384$    & $64$\,MB \\
\texttt{1024}  & $65\,536$    & $64$\,MB \\
\bottomrule
\end{tabular}
\end{center}

\begin{idea}
\textbf{Cuánto importa elegir \texttt{bs}}: bloques más grandes
significan menos syscalls (una por bloque). En la práctica, copiar
$64$\,MB con \texttt{bs=1} (un byte por vez) tarda decenas de
segundos, mientras que con \texttt{bs=1M} tarda fracciones de segundo.
La carga de trabajo no es la transferencia: son las llamadas al
sistema y los cambios de contexto.
\end{idea}

\subsection*{Solución}

La forma idiomática:

\begin{lstlisting}[style=shellstyle]
dd if=/dev/zero of=zero64.img bs=1M count=64
\end{lstlisting}

\textbf{Salida esperada.}

\begin{lstlisting}[style=textstyle]
64+0 records in
64+0 records out
67108864 bytes (67 MB, 64 MiB) copied, 0.046 s, 1.5 GB/s
\end{lstlisting}

\textbf{Verificación.}

\begin{lstlisting}[style=shellstyle]
$ ls -lh zero64.img
-rw-r--r-- 1 tomi tomi 64M ... zero64.img

$ wc -c zero64.img
67108864 zero64.img
\end{lstlisting}

\begin{importante}
\textbf{Alternativa moderna} para crear un archivo de $64$\,MB de
ceros: \verb|truncate -s 64M zero64.img|. La diferencia es que
\texttt{truncate} crea un \emph{sparse file}: el archivo aparenta tener
$64$\,MB pero no consume bloques en disco hasta que se escriben datos
no nulos. Con \texttt{dd} se escriben efectivamente $64$\,MB de ceros
físicos al disco.
\end{importante}

\textbf{Conclusión.} El comando solicitado es

\verb|dd if=/dev/zero of=zero64.img bs=1M count=64|.

% =====================================================================
\section{Ejercicio 3 --- Contraseña base64 desde \texttt{/dev/urandom}}
% =====================================================================

\textbf{Enunciado.} Escribir un comando que genere una contraseña de
$16$ caracteres codificada en base64 desde \texttt{/dev/urandom},
usando \texttt{head} y \texttt{base64}.

\subsection*{Análisis}

\textbf{Sobre \texttt{/dev/urandom}}: es un dispositivo de carácter
que entrega bytes pseudoaleatorios del PRNG del kernel. A diferencia
de \texttt{/dev/random}, \emph{no se bloquea} cuando el ``pool'' de
entropía se considera bajo; en sistemas modernos esto está
considerado seguro para cualquier uso criptográfico de propósito
general.

\textbf{Sobre base64}: codifica $3$ bytes de entrada en $4$
caracteres ASCII del alfabeto
$\{$A--Z, a--z, 0--9, +, /$\}$. La relación es:

\quad chars\_base64 = ⌈ bytes\_entrada × $4 \div 3$ ⌉

Despejando el camino inverso, para obtener \textbf{exactamente $16$
caracteres} sin padding ($=$):

\quad bytes\_entrada = $16 × 3 \div 4 = 12$

\subsection*{Solución}

\begin{lstlisting}[style=shellstyle]
head -c 12 /dev/urandom | base64
\end{lstlisting}

\begin{itemize}[leftmargin=*,nosep]
    \item \texttt{head -c 12 /dev/urandom}: lee exactamente $12$ bytes
        binarios del PRNG.
    \item Pipe a \texttt{base64}: codifica esos $12$ bytes en $16$
        caracteres ASCII.
\end{itemize}

\textbf{Ejemplos de salida.}

\begin{lstlisting}[style=textstyle]
$ head -c 12 /dev/urandom | base64
xK9q+Tb2/PfL3aRm

$ head -c 12 /dev/urandom | base64
hM7Bn5RpQ2Lk8VxC

$ head -c 12 /dev/urandom | base64
9j+kP/2nXqRtB4vL
\end{lstlisting}

\begin{idea}
Cada ejecución produce una contraseña distinta (es lo que se espera de
un generador aleatorio). Con $12$ bytes aleatorios bien distribuidos
hay $2^{12 \times 8} = 2^{96}$ posibilidades, mucho más que suficiente
para resistir cualquier ataque por fuerza bruta.
\end{idea}

\begin{importante}
\textbf{Para uso real} conviene asegurarse de eliminar el newline
final que \texttt{base64} agrega:

\begin{lstlisting}[style=shellstyle]
head -c 12 /dev/urandom | base64 | tr -d '\n'
\end{lstlisting}

\noindent Y para garantizar exactamente $16$ caracteres sin importar
padding:

\begin{lstlisting}[style=shellstyle]
head -c 16 /dev/urandom | base64 | head -c 16
\end{lstlisting}
\end{importante}

\textbf{Conclusión.} El comando pedido es

\verb|head -c 12 /dev/urandom | base64|.

% =====================================================================
\section{Ejercicio 4 --- Compilar y cargar el módulo RTC}
% =====================================================================

\textbf{Enunciado.} Compilar el módulo (device driver) Linux RTC
(\emph{Real Time Clock}) siguiendo las instrucciones del README. El
ejercicio se realiza sobre un sistema Linux con headers del kernel
instalados.

\subsection*{Análisis: qué es el RTC y por qué se accede como driver}

El \textbf{Real Time Clock} es un chip alimentado por batería que
mantiene la hora aunque la computadora esté apagada. En arquitectura
PC suele estar integrado al chipset (el clásico \texttt{MC146818}
emulado por los SO modernos). El kernel lo expone como dispositivo de
carácter en \texttt{/dev/rtc0}.

El RTC ofrece tres servicios al SO:

\begin{enumerate}[leftmargin=*,nosep]
    \item Lectura/escritura de la hora actual.
    \item Programación de \emph{alarmas} (interrupciones a futuro).
    \item Generación de \emph{interrupciones periódicas} (típicamente
        de $2$\,Hz a $8192$\,Hz), útiles para ticks de scheduler en
        sistemas embebidos.
\end{enumerate}

\subsection*{Anatomía mínima de un módulo del kernel}

Un módulo del kernel es código que se compila aparte del kernel
principal y se carga/descarga dinámicamente, sin reiniciar. El
esqueleto mínimo en C es:

\begin{lstlisting}[style=cstyle]
/* hello_rtc.c -- esqueleto de un modulo */
#include <linux/module.h>
#include <linux/kernel.h>
#include <linux/init.h>

static int __init hello_rtc_init(void) {
    printk(KERN_INFO "hello_rtc: modulo cargado\n");
    return 0;     /* 0 = exito, !=0 = falla la carga */
}

static void __exit hello_rtc_exit(void) {
    printk(KERN_INFO "hello_rtc: modulo descargado\n");
}

module_init(hello_rtc_init);
module_exit(hello_rtc_exit);
MODULE_LICENSE("GPL");
MODULE_AUTHOR("Curso SO UNRC");
MODULE_DESCRIPTION("Demo modulo RTC");
\end{lstlisting}

\textbf{Diferencias con un programa de usuario:}

\begin{itemize}[leftmargin=*,nosep]
    \item No hay \texttt{main()}: el ``punto de entrada'' es la función
        registrada con \texttt{module\_init}, que se llama cuando el
        módulo se carga.
    \item No se incluye \texttt{<stdio.h>}: en el kernel \texttt{printf}
        se reemplaza por \texttt{printk}.
    \item Hay un macro \texttt{KERN\_INFO} (y otros como
        \texttt{KERN\_WARNING}, \texttt{KERN\_ERR}) que marca el nivel
        de log del mensaje.
    \item El módulo se compila con un \texttt{Makefile} especial que
        usa el sistema de \emph{Kbuild} del kernel.
\end{itemize}

\subsection*{Makefile para compilar el módulo}

\begin{lstlisting}[style=textstyle]
# Makefile
obj-m += hello_rtc.o

KDIR := /lib/modules/$(shell uname -r)/build

all:
	$(MAKE) -C $(KDIR) M=$(PWD) modules

clean:
	$(MAKE) -C $(KDIR) M=$(PWD) clean
\end{lstlisting}

\begin{importante}
La línea \texttt{KDIR := /lib/modules/\$(shell uname -r)/build} pide
los \emph{headers} y scripts del kernel que se está ejecutando.
Necesita el paquete \texttt{linux-headers-\$(uname -r)} (Debian/Ubuntu)
o \texttt{kernel-devel} (Fedora/RHEL) instalado.
\end{importante}

\subsection*{Flujo de compilación, carga y prueba}

\begin{lstlisting}[style=shellstyle]
# 1) Compilar
$ make
make -C /lib/modules/6.5.0-15-generic/build M=/home/user/rtc modules
make[1]: Entering directory '/usr/src/linux-headers-6.5.0-15-generic'
  CC [M]  /home/user/rtc/hello_rtc.o
  MODPOST /home/user/rtc/Module.symvers
  CC [M]  /home/user/rtc/hello_rtc.mod.o
  LD [M]  /home/user/rtc/hello_rtc.ko

# 2) Cargar el modulo
$ sudo insmod hello_rtc.ko

# 3) Verificar que esta cargado
$ lsmod | grep hello_rtc
hello_rtc   16384  0

# 4) Ver el mensaje del modulo en el log del kernel
$ sudo dmesg | tail -n 2
[12345.678] hello_rtc: modulo cargado

# 5) Descargar
$ sudo rmmod hello_rtc

$ sudo dmesg | tail -n 1
[12346.789] hello_rtc: modulo descargado
\end{lstlisting}

\subsection*{Leer la hora del RTC desde el espacio de usuario}

Una vez cargado el driver del RTC (en general viene en el kernel por
defecto), se puede leer la hora con un programa de usuario que abra
\texttt{/dev/rtc0} y use \texttt{ioctl} con el comando
\texttt{RTC\_RD\_TIME}:

\begin{lstlisting}[style=cstyle]
/* read_rtc.c */
#include <stdio.h>
#include <fcntl.h>
#include <unistd.h>
#include <sys/ioctl.h>
#include <linux/rtc.h>

int main(void) {
    int fd = open("/dev/rtc0", O_RDONLY);
    if (fd < 0) { perror("open /dev/rtc0"); return 1; }

    struct rtc_time t;
    if (ioctl(fd, RTC_RD_TIME, &t) < 0) {
        perror("ioctl RTC_RD_TIME");
        close(fd);
        return 1;
    }

    printf("Hora del RTC: %04d-%02d-%02d %02d:%02d:%02d\n",
           t.tm_year + 1900, t.tm_mon + 1, t.tm_mday,
           t.tm_hour, t.tm_min, t.tm_sec);

    close(fd);
    return 0;
}
\end{lstlisting}

\begin{lstlisting}[style=shellstyle]
$ gcc -Wall -o read_rtc read_rtc.c
$ sudo ./read_rtc       # /dev/rtc0 generalmente requiere root
Hora del RTC: 2026-03-15 14:23:47
\end{lstlisting}

\subsection*{Conclusión}

El ejercicio ilustra el patrón estándar para extender el kernel con un
\textbf{driver de dispositivo} como módulo cargable:

\begin{enumerate}[leftmargin=*,nosep]
    \item Código C usando la API del kernel (\texttt{printk},
        \texttt{module\_init/exit}, etc.).
    \item Makefile que delega en Kbuild
        (\texttt{make -C /lib/modules/.../build M=\$PWD modules}).
    \item Carga/descarga con \texttt{insmod}/\texttt{rmmod}.
    \item Inspección con \texttt{lsmod}, \texttt{dmesg},
        \texttt{/proc/modules}.
\end{enumerate}

El RTC, en particular, es accedido desde el espacio de usuario
mediante \texttt{ioctl} sobre \texttt{/dev/rtc0}, ejemplo paradigmático
de cómo el SO ofrece una interfaz uniforme (un archivo) para un
dispositivo que internamente requiere comandos específicos.

% =====================================================================
\section{Ejercicio 5 --- Filesystem \texttt{/docs/os-book.pdf}: FAT vs.\ inodos}
% =====================================================================

\textbf{Enunciado.} Describir un sistema de archivos con la estructura

\begin{lstlisting}[style=textstyle]
\
|-- docs
       |
       +-- os-book.pdf      (2 bloques de datos)
\end{lstlisting}

\noindent usando \textbf{(a)} una representación tipo FAT y
\textbf{(b)} una representación basada en \emph{inodos}.

\subsection*{Análisis comparativo}

Las dos familias de filesystem responden de forma distinta a la
pregunta ``dado el nombre de un archivo, ¿cómo encuentro sus
bloques?''.

\begin{itemize}[leftmargin=*,nosep]
    \item En \textbf{FAT}, la entrada de directorio guarda el
        \emph{primer cluster} del archivo. Para los demás, se consulta
        una \textbf{tabla centralizada} (la FAT) que es esencialmente
        una lista enlazada: la entrada \texttt{FAT[i]} dice cuál es el
        próximo cluster del archivo que comenzó en $i$, o un valor
        especial $\texttt{EOC}$ si $i$ es el último.
    \item En \textbf{inodos} (estilo UFS/ext2/\ldots), la entrada de
        directorio guarda sólo el número de \emph{inodo}. El inodo es
        un descriptor con toda la metadata del archivo (tamaño,
        permisos, owner, fechas) y \textbf{punteros directos} a sus
        bloques. La estructura está \emph{descentralizada}: cada
        archivo lleva su propia lista.
\end{itemize}

\subsection*{5.a --- Representación tipo FAT}

Una partición FAT tiene cuatro áreas principales: \emph{boot sector},
\emph{tabla FAT} (a veces duplicada), \emph{directorio raíz} y
\emph{área de datos} (clusters numerados a partir de $2$, porque $0$ y
$1$ están reservados).

\paragraph{Entradas de directorio.} Cada entrada es de tamaño fijo
(usualmente $32$ bytes) y contiene:

\begin{center}
\small
\begin{tabular}{@{}lp{8.5cm}@{}}
\toprule
Campo & Significado \\
\midrule
\texttt{name}        & nombre $8 + 3$ chars (FAT16) o de tamaño variable (FAT32 con LFN) \\
\texttt{attr}        & directorio, archivo, oculto, sólo lectura, \ldots \\
\texttt{first\_clu}  & primer cluster del archivo en el área de datos \\
\texttt{size}        & tamaño en bytes (sólo para archivos regulares) \\
\bottomrule
\end{tabular}
\end{center}

\paragraph{Layout concreto del ejemplo.} Supongamos:

\begin{itemize}[leftmargin=*,nosep]
    \item Cluster $3$: contenido del directorio \texttt{docs}.
    \item Cluster $5$: primer bloque del PDF.
    \item Cluster $6$: segundo bloque del PDF.
\end{itemize}

\textbf{Directorio raíz} (área especial en FAT16, o cluster fijo en
FAT32). Una sola entrada:

\begin{center}
\small
\begin{tabular}{@{}llll@{}}
\toprule
name & attr & first\_clu & size \\
\midrule
\texttt{DOCS} & DIR & $3$ & $0$ \\
\bottomrule
\end{tabular}
\end{center}

\textbf{Contenido del cluster $3$} (directorio \texttt{docs}):

\begin{center}
\small
\begin{tabular}{@{}llll@{}}
\toprule
name & attr & first\_clu & size \\
\midrule
\texttt{.}            & DIR & $3$ & $0$       \\
\texttt{..}           & DIR & $0$ (raíz)  & $0$       \\
\texttt{OS-BOOK PDF}  & --- & $5$ & $2 \times $bs \\
\bottomrule
\end{tabular}
\end{center}

\textbf{Tabla FAT.} Cada entrada \texttt{FAT[i]} dice quién sigue al
cluster $i$:

\begin{center}
\small
\begin{tabular}{@{}cl@{}}
\toprule
i & FAT[i] \\
\midrule
$0$ & reservado (marcador) \\
$1$ & reservado \\
$2$ & libre ($0$) \\
$3$ (raíz de \texttt{docs}) & \texttt{EOC} \\
$4$ & libre \\
$5$ (\texttt{os-book.pdf} parte $1$) & $6$ \\
$6$ (\texttt{os-book.pdf} parte $2$) & \texttt{EOC} \\
$\geq 7$ & libre \\
\bottomrule
\end{tabular}
\end{center}

\begin{idea}
Para leer \texttt{/docs/os-book.pdf} el SO recorre:

\begin{enumerate}[leftmargin=*,nosep,label=(\arabic*)]
    \item Raíz: encuentra entrada \texttt{DOCS}, \texttt{first\_clu = 3}.
    \item Va al cluster $3$, encuentra entrada
        \texttt{OS-BOOK.PDF}, \texttt{first\_clu = 5}.
    \item Lee el cluster $5$.
    \item Consulta \texttt{FAT[5]} = $6$ $\to$ lee el cluster $6$.
    \item Consulta \texttt{FAT[6]} = \texttt{EOC} $\to$ termina.
\end{enumerate}
\end{idea}

\subsection*{Diagrama del layout FAT}

\begin{center}
\begin{tikzpicture}[
    every node/.style={font=\ttfamily\footnotesize},
    sec/.style={draw,minimum width=3.2cm,minimum height=0.7cm,align=center},
    cl/.style={draw,minimum width=1.8cm,minimum height=0.6cm,align=center}
]

\node[font=\bfseries\small] at (0,5.5) {Disco};
\node[sec,fill=io1!20] (boot) at (0,4.7) {boot sector};
\node[sec,fill=io2!20,below=0pt of boot] (fat) {tabla FAT};
\node[sec,fill=io3!20,below=0pt of fat] (rd) {root dir};
\node[sec,fill=ioshared!30,below=0pt of rd] (data) {\ldots área de datos\ldots};

\node[font=\bfseries\small] at (5,5.5) {Clusters};
\node[cl,fill=iofree!30] (c2) at (5,5) {2 (libre)};
\node[cl,fill=io3!20,below=0pt of c2] (c3) {3: dir docs};
\node[cl,fill=iofree!30,below=0pt of c3] (c4) {4 (libre)};
\node[cl,fill=iofile!50,below=0pt of c4] (c5) {5: pdf p1};
\node[cl,fill=iofile!50,below=0pt of c5] (c6) {6: pdf p2};
\node[cl,fill=iofree!30,below=0pt of c6] (c7) {\ldots};

\node[font=\bfseries\small] at (10,5.5) {FAT[i]};
\node[cl,fill=io2!10] (f2) at (10,5) {2 $\to$ 0};
\node[cl,fill=io2!10,below=0pt of f2] (f3) {3 $\to$ EOC};
\node[cl,fill=io2!10,below=0pt of f3] (f4) {4 $\to$ 0};
\node[cl,fill=io2!10,below=0pt of f4] (f5) {5 $\to$ 6};
\node[cl,fill=io2!10,below=0pt of f5] (f6) {6 $\to$ EOC};
\node[cl,fill=io2!10,below=0pt of f6] (f7) {\ldots};

\draw[-Latex] (c5.east) to[bend left=15] (c6.east);
\draw[-Latex,dashed,thick,red!60!black] (f5.west) -- (c5.east);
\draw[-Latex,dashed,thick,red!60!black] (f6.west) -- (c6.east);

\end{tikzpicture}
\end{center}

\subsection*{5.b --- Representación con inodos}

En un FS basado en inodos cada archivo o directorio tiene \emph{un}
inodo numerado de forma única dentro del FS. El inodo guarda toda la
metadata y los punteros a bloques de datos. Las entradas de directorio
sólo asocian nombres con números de inodo.

\paragraph{Estructura típica de un inodo} (estilo UFS clásico):

\begin{center}
\small
\begin{tabular}{@{}lp{8cm}@{}}
\toprule
Campo & Significado \\
\midrule
\texttt{i\_mode}                 & tipo (regular, dir, link, \ldots) y permisos \\
\texttt{i\_uid}, \texttt{i\_gid} & owner y grupo \\
\texttt{i\_size}                 & tamaño en bytes \\
\texttt{i\_atime, mtime, ctime}  & timestamps de acceso, modificación, cambio de inodo \\
\texttt{i\_blocks}               & bloques físicos usados \\
\texttt{i\_direct[0..11]}        & punteros directos a bloques de datos \\
\texttt{i\_indirect1}            & puntero indirecto simple \\
\texttt{i\_indirect2}            & puntero indirecto doble \\
\texttt{i\_indirect3}            & puntero indirecto triple \\
\bottomrule
\end{tabular}
\end{center}

\paragraph{Asignación de inodos al ejemplo.}

\begin{center}
\small
\begin{tabular}{@{}clp{6.5cm}@{}}
\toprule
\# inodo & Quién es & Apunta a (bloques de datos) \\
\midrule
$1$ & \texttt{/} (raíz, directorio)        & $B_{root}$ \\
$2$ & \texttt{/docs} (directorio)           & $B_{docs}$ \\
$3$ & \texttt{/docs/os-book.pdf} (regular) & \texttt{direct[0]=$B_{pdf1}$}, \texttt{direct[1]=$B_{pdf2}$} \\
\bottomrule
\end{tabular}
\end{center}

\paragraph{Contenido de los bloques de directorio} (cada entrada =
nombre $\to$ inodo):

\begin{itemize}[leftmargin=*,nosep]
    \item $B_{root}$:
        $\{\;\texttt{.} \to 1,\;\; \texttt{..} \to 1,\;\; \texttt{docs} \to 2\;\}$
    \item $B_{docs}$:
        $\{\;\texttt{.} \to 2,\;\; \texttt{..} \to 1,\;\; \texttt{os-book.pdf} \to 3\;\}$
    \item $B_{pdf1}$, $B_{pdf2}$: bytes binarios del PDF.
\end{itemize}

\textbf{Notas.}

\begin{itemize}[leftmargin=*,nosep]
    \item La raíz tiene siempre el mismo número de inodo (en ext2/3/4
        el $\#2$, no el $\#1$, que se usa para tracking de bloques
        dañados; el $\#1$ usado aquí es ilustrativo).
    \item \texttt{..} de la raíz apunta a la raíz misma.
    \item \texttt{i\_size} del inodo $3$ es $\leq 2 × \text{bs}$ (no
        necesariamente exactamente $2$ bloques: el archivo puede
        terminar a mitad del segundo).
\end{itemize}

\subsection*{Diagrama del layout con inodos}

\begin{center}
\begin{tikzpicture}[
    every node/.style={font=\ttfamily\footnotesize},
    inode/.style={draw,minimum width=3.6cm,minimum height=0.55cm,align=center},
    block/.style={draw,minimum width=3.6cm,minimum height=0.55cm,align=center}
]

\node[font=\bfseries\small] at (0,5.5) {Tabla de inodos};
\node[inode,fill=io1!20] (i1) at (0,4.8) {\#1: DIR, $B_{root}$};
\node[inode,fill=io1!20,below=0pt of i1] (i2) {\#2: DIR, $B_{docs}$};
\node[inode,fill=io2!20,below=0pt of i2] (i3) {\#3: REG, $B_{pdf1}$, $B_{pdf2}$};

\node[font=\bfseries\small] at (6,5.5) {Bloques de datos};
\node[block,fill=io3!15] (br) at (6,4.8) {$B_{root}$: \{., .., docs\}};
\node[block,fill=io3!15,below=0pt of br] (bd) {$B_{docs}$: \{., .., os-book.pdf\}};
\node[block,fill=iofile!50,below=0pt of bd] (bp1) {$B_{pdf1}$: bytes PDF (1/2)};
\node[block,fill=iofile!50,below=0pt of bp1] (bp2) {$B_{pdf2}$: bytes PDF (2/2)};

\draw[-Latex] (i1.east) -- (br.west);
\draw[-Latex] (i2.east) -- (bd.west);
\draw[-Latex] (i3.east) -- (bp1.west);
\draw[-Latex] (i3.east) to[bend right=15] (bp2.west);

\end{tikzpicture}
\end{center}

\subsection*{Conclusión: contraste FAT vs.\ inodos}

\begin{center}
\small
\begin{tabular}{@{}lll@{}}
\toprule
 & \textbf{FAT} & \textbf{Inodos} \\
\midrule
Metadata del archivo & en la entrada de directorio  & en el inodo \\
Lista de bloques     & FAT centralizada (lista enlazada) & punteros en el inodo \\
\emph{Hard links} (varios nombres) & \textbf{No} soportado & Sí: varias entradas $\to$ mismo inodo \\
Acceso aleatorio rápido & no (recorrer FAT)        & sí (punteros directos / indirectos) \\
Robustez ante FAT corrupta & toda la partición se pierde & sólo se afecta un archivo \\
\bottomrule
\end{tabular}
\end{center}

\textbf{Resumen.} La diferencia esencial es \emph{dónde} se guarda la
información de qué bloques pertenecen a qué archivo. En FAT, toda la
información de pertenencia está en una sola tabla compartida. En
inodos, cada archivo lleva su propia ``tabla'' (los punteros del
inodo), y los directorios se reducen a una asociación de nombres con
números de inodo.

% =====================================================================
\section{Ejercicio 6 --- Mostrar tipo y tamaño con \texttt{stat()}}
% =====================================================================

\textbf{Enunciado.} El siguiente programa lee el directorio corriente
e imprime el nombre y el inodo de cada entrada. Extenderlo para
mostrar también el \textbf{tipo} (regular, directorio, link, \ldots)
y el \textbf{tamaño} de cada archivo, usando \texttt{stat()} sobre
\texttt{st\_mode} y \texttt{st\_size}.

\subsection*{Análisis: por qué necesitamos \texttt{stat()}}

\texttt{readdir()} devuelve una \texttt{struct dirent} con el
\emph{nombre} y el \emph{número de inodo} de la entrada, pero \textbf{no}
con la metadata completa (tipo, tamaño, permisos, fechas). Esa
información vive en el inodo, no en la entrada de directorio.

Para obtenerla hay que invocar \texttt{stat()} (o sus variantes
\texttt{lstat}, \texttt{fstat}), que consulta el inodo del archivo y
llena una \texttt{struct stat} con todos los campos:

\begin{lstlisting}[style=cstyle]
struct stat {
    dev_t     st_dev;       /* dispositivo */
    ino_t     st_ino;       /* numero de inodo */
    mode_t    st_mode;      /* tipo + permisos */
    nlink_t   st_nlink;     /* cantidad de hard links */
    uid_t     st_uid;       /* owner */
    gid_t     st_gid;       /* grupo */
    off_t     st_size;      /* tamano en bytes */
    blksize_t st_blksize;   /* tamano del bloque */
    blkcnt_t  st_blocks;    /* bloques fisicos asignados */
    time_t    st_atime, st_mtime, st_ctime;
    /* ... */
};
\end{lstlisting}

\textbf{Tipo del archivo en \texttt{st\_mode}.} El campo combina dos
informaciones: tipo y permisos. El tipo se extrae con macros
\texttt{S\_IS<Tipo>(m)} provistas por \texttt{<sys/stat.h>}:

\begin{center}
\small
\begin{tabular}{@{}lp{8cm}@{}}
\toprule
Macro              & Devuelve no-cero si\ldots \\
\midrule
\texttt{S\_ISREG(m)}  & el archivo es \textbf{regular} \\
\texttt{S\_ISDIR(m)}  & el archivo es un \textbf{directorio} \\
\texttt{S\_ISLNK(m)}  & el archivo es un \textbf{symlink} \\
\texttt{S\_ISCHR(m)}  & el archivo es un \textbf{character device} \\
\texttt{S\_ISBLK(m)}  & el archivo es un \textbf{block device} \\
\texttt{S\_ISFIFO(m)} & el archivo es una \textbf{FIFO / pipe con nombre} \\
\texttt{S\_ISSOCK(m)} & el archivo es un \textbf{socket de UNIX} \\
\bottomrule
\end{tabular}
\end{center}

\begin{importante}
\textbf{\texttt{stat} vs \texttt{lstat}.} Si la entrada de directorio
es un \textbf{symlink}, \texttt{stat()} sigue el link y reporta los
datos del \emph{archivo apuntado}, no del link mismo. Para obtener los
datos del link en sí, hay que usar \texttt{lstat()}. En este programa
conviene \texttt{lstat()} para no confundir symlinks con sus targets.
\end{importante}

\subsection*{Solución}

\begin{lstlisting}[style=cstyle]
#include <stdio.h>
#include <string.h>
#include <dirent.h>
#include <sys/stat.h>

static const char *file_type(mode_t m) {
    if (S_ISREG(m))  return "REG ";
    if (S_ISDIR(m))  return "DIR ";
    if (S_ISLNK(m))  return "LNK ";
    if (S_ISCHR(m))  return "CHR ";
    if (S_ISBLK(m))  return "BLK ";
    if (S_ISFIFO(m)) return "FIFO";
    if (S_ISSOCK(m)) return "SOCK";
    return "?   ";
}

int main(void) {
    DIR *dir = opendir(".");
    if (!dir) { perror("opendir"); return 1; }

    struct dirent *entry;
    while ((entry = readdir(dir)) != NULL) {
        struct stat st;
        if (lstat(entry->d_name, &st) < 0) {
            perror(entry->d_name);
            continue;
        }
        printf("%-25s  inode=%-10lu  type=%s  size=%lld bytes\n",
               entry->d_name,
               (unsigned long)entry->d_ino,
               file_type(st.st_mode),
               (long long)st.st_size);
    }
    closedir(dir);
    return 0;
}
\end{lstlisting}

\subsection*{Salida típica}

Ejecutándolo en un directorio con un poco de variedad:

\begin{lstlisting}[style=textstyle]
$ gcc -Wall -o lsdir lsdir.c
$ ./lsdir
.                          inode=14385       type=DIR   size=4096 bytes
..                         inode=14112       type=DIR   size=4096 bytes
lsdir.c                    inode=14502       type=REG   size=812 bytes
lsdir                      inode=14503       type=REG   size=16400 bytes
notas.txt                  inode=14504       type=REG   size=125 bytes
backup -> notas.txt        inode=14505       type=LNK   size=10 bytes
\end{lstlisting}

\subsection*{Lectura de los resultados}

\begin{itemize}[leftmargin=*]
    \item \textbf{Directorios} (\texttt{.}, \texttt{..}): el campo
        \texttt{size} es el tamaño \emph{del bloque de directorio} (en
        ext4 típicamente $4096$\,B), NO la suma del contenido. Para
        saber cuánto pesan los archivos de adentro hay que recorrer
        el directorio recursivamente (lo que hace \texttt{du}).
    \item \textbf{Archivos regulares}: \texttt{size} es el tamaño en
        bytes del contenido.
    \item \textbf{Symlinks}: \texttt{size} con \texttt{lstat} es la
        \emph{longitud del target string} (en este caso $10$, los
        chars de \texttt{"notas.txt"}). Con \texttt{stat}, en cambio,
        habría dado $125$, el tamaño del archivo apuntado.
\end{itemize}

\textbf{Conclusión.} Las dos pistas claves para resolver el ejercicio
fueron: (1) usar \texttt{stat}/\texttt{lstat} para obtener metadata
adicional sobre cada entrada de directorio, y (2) decodificar
\texttt{st\_mode} con las macros \texttt{S\_IS<Tipo>(m)} para mapear el
bitfield a una etiqueta legible.

% =====================================================================
\part*{Bloque B --- Usuarios y control de acceso}
\addcontentsline{toc}{section}{Bloque B --- Usuarios y control de acceso}
% =====================================================================

% =====================================================================
\section{Ejercicio 7 --- Listar usuarios y grupos del sistema}
% =====================================================================

\textbf{Enunciado.} Listar los usuarios y grupos del sistema. En macOS,
usar el comando \texttt{dscl} (\emph{directory services command-line}).

\subsection*{Análisis}

Un sistema UNIX mantiene la información de cuentas en dos archivos
históricos:

\begin{itemize}[leftmargin=*,nosep]
    \item \texttt{/etc/passwd}: una línea por usuario, $7$ campos
        separados por \texttt{:}.
    \item \texttt{/etc/group}: una línea por grupo, $4$ campos.
\end{itemize}

En sistemas modernos puede no usarse este archivo directamente (por
ejemplo, cuando hay LDAP, Active Directory u Open Directory en macOS),
pero la API portable sigue siendo la misma: \texttt{getpwent},
\texttt{getgrent} en C, o \texttt{getent} en shell.

\subsection*{Linux}

\textbf{Usuarios.}

\begin{lstlisting}[style=shellstyle]
$ cat /etc/passwd | head -n 4
root:x:0:0:root:/root:/bin/bash
daemon:x:1:1:daemon:/usr/sbin:/usr/sbin/nologin
bin:x:2:2:bin:/bin:/usr/sbin/nologin
sys:x:3:3:sys:/dev:/usr/sbin/nologin

# Forma portable: consulta la base de datos del sistema
$ getent passwd | wc -l
38
\end{lstlisting}

\textbf{Formato de \texttt{/etc/passwd}.}

\begin{center}
\small
\begin{tabular}{@{}clp{8cm}@{}}
\toprule
\# & Campo & Significado \\
\midrule
$1$ & \texttt{name}   & login del usuario \\
$2$ & \texttt{passwd} & contraseña (\texttt{x} significa ``en \texttt{/etc/shadow}'') \\
$3$ & \texttt{uid}    & identificador numérico \\
$4$ & \texttt{gid}    & grupo primario \\
$5$ & \texttt{gecos}  & información extendida (nombre real, etc.) \\
$6$ & \texttt{home}   & directorio personal \\
$7$ & \texttt{shell}  & shell de login (\texttt{/bin/bash}, \texttt{/usr/sbin/nologin}, \ldots) \\
\bottomrule
\end{tabular}
\end{center}

\textbf{Grupos.}

\begin{lstlisting}[style=shellstyle]
$ cat /etc/group | head -n 4
root:x:0:
daemon:x:1:
bin:x:2:
sys:x:3:

$ getent group sudo
sudo:x:27:tomi,otra_persona
\end{lstlisting}

\textbf{Formato de \texttt{/etc/group}.} \texttt{name : passwd : gid :
members\_separados\_por\_coma}.

\subsection*{macOS}

macOS no usa \texttt{/etc/passwd} para usuarios reales (sólo para
algunas cuentas del sistema). La fuente de verdad es \textbf{Open
Directory}, accesible vía \texttt{dscl}:

\begin{lstlisting}[style=shellstyle]
# Listar todos los usuarios
$ dscl . -list /Users
_amavisd
_appleevents
_appowner
...
root
tomi

# Filtrar solo los "humanos" (los del sistema empiezan con _)
$ dscl . -list /Users | grep -v '^_'
daemon
nobody
root
tomi

# Listar todos los grupos
$ dscl . -list /Groups | head
_amavisd
_appleevents
...

# Detalle de un usuario
$ dscl . -read /Users/tomi
NFSHomeDirectory: /Users/tomi
PrimaryGroupID: 20
RealName: Tomas Rodeghiero
RecordName: tomi
UniqueID: 501
UserShell: /bin/zsh
\end{lstlisting}

\begin{idea}
\textbf{Convención macOS}: los usuarios y grupos del \emph{sistema}
(no destinados a humanos) empiezan con \texttt{\_}. Filtrar con
\texttt{grep -v '\textasciicircum\_'} muestra sólo cuentas regulares.
\end{idea}

\textbf{Conclusión.} En Linux la forma portable es
\texttt{getent passwd}/\texttt{getent group} (consulta cualquier
backend configurado en \texttt{/etc/nsswitch.conf}); en macOS, el
comando equivalente es \texttt{dscl . -list /Users} y
\texttt{dscl . -list /Groups}.

% =====================================================================
\section{Ejercicio 8 --- Permisos UGO y octal sobre archivos y directorios}
% =====================================================================

\textbf{Enunciado.} Crear un directorio dentro del directorio de
trabajo. Analizar los permisos asignados, qué significa el bit de
ejecución sobre un directorio, crear un archivo dentro con
\texttt{touch}, manipular sus permisos con UGO y octal, y observar el
efecto.

\subsection*{Análisis: permisos en UNIX}

Cada archivo tiene una tripla de permisos:

\begin{center}
\small
\begin{tabular}{@{}lll@{}}
\toprule
Categoría & Letra & Quién \\
\midrule
Owner    & \texttt{u} & el dueño del archivo \\
Group    & \texttt{g} & los miembros del grupo del archivo \\
Other    & \texttt{o} & cualquier otro usuario \\
All      & \texttt{a} & \texttt{u + g + o} (atajo) \\
\bottomrule
\end{tabular}
\end{center}

Y tres bits por cada categoría:

\begin{center}
\small
\begin{tabular}{@{}lll@{}}
\toprule
Bit & Sobre archivo & Sobre directorio \\
\midrule
\texttt{r} ($4$) & leer contenido            & listar nombres (\texttt{ls}) \\
\texttt{w} ($2$) & modificar contenido       & crear/borrar/renombrar entradas \\
\texttt{x} ($1$) & ejecutar como programa    & atravesar (\texttt{cd}, acceder a entradas por nombre) \\
\bottomrule
\end{tabular}
\end{center}

\subsection*{8.a --- Permisos por defecto y el bit \texttt{x} en directorios}

\begin{lstlisting}[style=shellstyle]
$ mkdir mydir
$ ls -ld mydir
drwxr-xr-x  2 tomi staff  64 ... mydir
\end{lstlisting}

\textbf{Lectura de \texttt{drwxr-xr-x}.}

\begin{center}
\small
\begin{tabular}{@{}lllll@{}}
\toprule
Posición & \texttt{d} & \texttt{rwx} & \texttt{r-x} & \texttt{r-x} \\
Significado & directorio & owner: lee, escribe, atraviesa & group: lee, atraviesa & other: lee, atraviesa \\
Octal & --- & $7$ & $5$ & $5$ \\
\bottomrule
\end{tabular}
\end{center}

Es decir, $755$. ¿Por qué $755$ y no $777$? Por el \emph{umask},
una máscara de bits que el shell aplica al permiso ``base''
($777$ para directorios). El umask típico es $022$:

\quad permiso\_final = $0777$ \texttt{AND NOT} umask = $0777$ \texttt{AND NOT} $022$ = $755$

\paragraph{¿Qué significa el bit \texttt{x} en un directorio?}

\begin{itemize}[leftmargin=*]
    \item \textbf{Con} \texttt{x}: se puede ``entrar'' al directorio
        con \texttt{cd}, abrir archivos por su nombre y atravesarlo
        como parte de una ruta más larga.
    \item \textbf{Sin} \texttt{x}: aun teniendo \texttt{r} se pueden
        \emph{listar nombres} (\texttt{ls mydir}) pero \emph{no acceder
        al contenido} (\texttt{ls -l} falla porque necesita leer el
        inodo de cada entrada, lo cual requiere atravesar).
    \item \textbf{Con \texttt{r} y sin \texttt{x}}: ver nombres pero
        no acceder.
    \item \textbf{Con \texttt{x} y sin \texttt{r}}: acceder a archivos
        cuyos nombres se conozcan de antemano, pero no listar.
\end{itemize}

\begin{idea}
La pareja \texttt{r}/\texttt{x} sobre directorios es notoria por
confundir: \texttt{r} es ``ver la guía telefónica'' (listar entradas)
y \texttt{x} es ``llamar a un número conocido'' (acceder a un
archivo). Son ortogonales.
\end{idea}

\subsection*{8.b --- Crear un archivo con \texttt{touch} y mirar permisos}

\begin{lstlisting}[style=shellstyle]
$ cd mydir
$ touch test.txt
$ ls -l test.txt
-rw-r--r--  1 tomi staff  0 ... test.txt
\end{lstlisting}

Permisos = $644$.

\quad permiso\_final = $0666$ \texttt{AND NOT} $022$ = $0644$

(Para archivos regulares la base es $666$, no $777$: \textbf{los
archivos no nacen con \texttt{x}}; hay que ponerlo explícitamente con
\texttt{chmod +x}.)

\subsection*{8.c --- Quitar el permiso de lectura}

\textbf{Sintaxis UGO} (modificación simbólica relativa):

\begin{lstlisting}[style=shellstyle]
$ chmod a-r test.txt        # 'a' = u+g+o, '-r' quita el bit r
$ ls -l test.txt
--w-------  1 tomi staff  0 ... test.txt
\end{lstlisting}

\textbf{Sintaxis octal} (especificación absoluta):

\begin{lstlisting}[style=shellstyle]
$ chmod 200 test.txt        # 200 = u=w, g=---, o=---
$ ls -l test.txt
--w-------  1 tomi staff  0 ... test.txt
\end{lstlisting}

Las dos son equivalentes. UGO es más explícita en \emph{qué cambia};
octal es más explícita en \emph{cómo queda al final}.

\begin{center}
\small
\begin{tabular}{@{}cccc@{}}
\toprule
Octal & Binario & UGO & Permiso \\
\midrule
$7$ & $111$ & \texttt{rwx} & lee + escribe + ejecuta \\
$6$ & $110$ & \texttt{rw-} & lee + escribe \\
$5$ & $101$ & \texttt{r-x} & lee + ejecuta \\
$4$ & $100$ & \texttt{r--} & lee \\
$3$ & $011$ & \texttt{-wx} & escribe + ejecuta \\
$2$ & $010$ & \texttt{-w-} & escribe \\
$1$ & $001$ & \texttt{--x} & ejecuta \\
$0$ & $000$ & \texttt{---} & nada \\
\bottomrule
\end{tabular}
\end{center}

\subsection*{8.d --- Intentar editar el archivo}

Con $200$ el dueño tiene sólo \texttt{w}. Veamos qué pasa:

\begin{lstlisting}[style=shellstyle]
# Leer falla
$ cat test.txt
cat: test.txt: Permission denied

# Editar con vim/nano: depende del editor
$ vim test.txt
"test.txt" Permission denied
Press ENTER or type command to continue

# Sobrescribir con redireccion: FUNCIONA, porque solo necesita 'w'
$ echo "hola" > test.txt

# Append: tambien funciona (solo 'w', no 'r')
$ echo "otra linea" >> test.txt
\end{lstlisting}

\begin{importante}
\textbf{Editar con un editor no es lo mismo que escribir.} Los
editores como \texttt{vim}, \texttt{nano} o \texttt{emacs} necesitan
\emph{leer} el contenido actual antes de mostrarlo. Sin \texttt{r},
fallan al abrirlo. Pero el shell puede sobrescribir un archivo con
\texttt{>} usando solo el bit \texttt{w}. Es una distinción importante:
``no poder editar'' es más restrictivo que ``no poder leer''.
\end{importante}

\subsection*{8.e --- Restaurar la lectura sólo al dueño}

\textbf{UGO:}

\begin{lstlisting}[style=shellstyle]
$ chmod u+r test.txt        # agrega r solo al owner
$ ls -l test.txt
-rw-------  1 tomi staff  ...
\end{lstlisting}

\textbf{Octal equivalente:}

\begin{lstlisting}[style=shellstyle]
$ chmod 600 test.txt
\end{lstlisting}

\subsection*{8.f --- Volver a editar}

Con $600$ el dueño tiene \texttt{rw}, así que cualquier editor abre el
archivo sin problemas:

\begin{lstlisting}[style=shellstyle]
$ vim test.txt    # abre normalmente, edita, guarda
\end{lstlisting}

\textbf{Conclusión.} Los puntos clave del ejercicio:

\begin{enumerate}[leftmargin=*]
    \item Los permisos por defecto los fija el \texttt{umask}
        (típicamente $022$): $755$ para directorios, $644$ para
        archivos.
    \item El bit \texttt{x} en un directorio es ``permiso para
        atravesar'', no ``permiso para ejecutar''.
    \item UGO y octal son dos formas de la misma operación; cada bit
        corresponde a una posición en el octal ($r=4$, $w=2$, $x=1$).
    \item Sin \texttt{r} se puede aún \emph{sobrescribir} con \texttt{>}
        pero no abrir con un editor; los editores necesitan ambas
        cosas.
\end{enumerate}

% =====================================================================
\section{Ejercicio 9 --- Real UID, Effective UID y shell con \texttt{getpwuid}}
% =====================================================================

\textbf{Enunciado.} Sobre el programa dado, que imprime el real UID y
el effective UID con su nombre asociado vía \texttt{getpwuid()},
\textbf{extenderlo} para que también muestre el \emph{shell} que usa
el real user.

\subsection*{Análisis: real UID vs.\ effective UID}

Todo proceso UNIX tiene cuatro identificadores de usuario asociados:

\begin{center}
\small
\begin{tabular}{@{}lp{9cm}@{}}
\toprule
ID & Significado \\
\midrule
\texttt{ruid} (real)            & quién \textbf{lanzó} el proceso. Casi siempre fijo durante la vida del proceso. \\
\texttt{euid} (effective)       & quién es ``a los efectos de chequear permisos'' en cada syscall. \\
\texttt{suid} (saved-set)       & valor previo del \texttt{euid}; permite swap entre ruid/euid. \\
\texttt{fsuid} (filesystem)     & sólo Linux: \texttt{euid} a efectos de operaciones de FS. \\
\bottomrule
\end{tabular}
\end{center}

En la mayoría de los procesos \texttt{ruid = euid = suid}. La
distinción se vuelve relevante con programas SUID/SGID (Ejercicio 10).

\textbf{\texttt{getpwuid(uid)}} consulta la base de datos de usuarios
(\texttt{/etc/passwd} o equivalente) por el UID dado y devuelve una
\texttt{struct passwd *}:

\begin{lstlisting}[style=cstyle]
struct passwd {
    char   *pw_name;     /* login name        */
    char   *pw_passwd;   /* hashed password   */
    uid_t   pw_uid;      /* user id           */
    gid_t   pw_gid;      /* group id          */
    char   *pw_gecos;    /* descripcion       */
    char   *pw_dir;      /* home directory    */
    char   *pw_shell;    /* default shell     */
};
\end{lstlisting}

\subsection*{Solución}

\begin{lstlisting}[style=cstyle]
#include <stdio.h>
#include <unistd.h>
#include <sys/types.h>
#include <pwd.h>

int main(void) {
    uid_t ruid = getuid();   /* real UID */
    uid_t euid = geteuid();  /* effective UID */

    struct passwd *pw_r = getpwuid(ruid);
    if (pw_r == NULL) { perror("getpwuid(ruid)"); return 1; }

    printf("Process Real User ID: %d, name: %s\n", ruid, pw_r->pw_name);
    printf("Real user shell     : %s\n", pw_r->pw_shell);    /* <-- agregado */

    struct passwd *pw_e = getpwuid(euid);
    if (pw_e == NULL) { perror("getpwuid(euid)"); return 1; }

    printf("Process Effective User ID: %d, name: %s\n", euid, pw_e->pw_name);

    return 0;
}
\end{lstlisting}

\subsection*{Compilación y ejecución}

\begin{lstlisting}[style=shellstyle]
$ gcc -Wall -o whoami_extended whoami_extended.c
$ ./whoami_extended
Process Real User ID: 1000, name: tomi
Real user shell     : /bin/bash
Process Effective User ID: 1000, name: tomi
\end{lstlisting}

\textbf{Observación.} En una ejecución normal \texttt{ruid = euid}.
La diferencia aparecerá en el Ejercicio 10 cuando se active el bit
SUID.

\begin{idea}
El campo \texttt{pw\_shell} es el \emph{login shell} configurado para
el usuario (séptimo campo de \texttt{/etc/passwd}). No siempre es el
shell que el usuario está usando ahora: si abrió una sub-sesión con
\texttt{zsh}, el valor de \texttt{pw\_shell} sigue mostrando el shell
por defecto. Para conocer el shell en uso \emph{en este preciso
proceso} hay que consultar la variable de ambiente \texttt{\$SHELL}
del entorno heredado.
\end{idea}

\textbf{Conclusión.} La extensión es de una sola línea: imprimir
\texttt{pw\_r->pw\_shell} después de obtener el \texttt{struct passwd}
del real UID. Pero el ejercicio fuerza a entender la diferencia entre
el UID \emph{real} (quién lanzó el proceso) y el \emph{efectivo} (con
qué identidad se evalúan los permisos en cada syscall).

% =====================================================================
\section{Ejercicio 10 --- El bit SUID en acción}
% =====================================================================

\textbf{Enunciado.} Extender el programa anterior para que (a) abra
el archivo creado en el ejercicio 8 para lectura/escritura, (b) se
cambie el dueño del archivo a \texttt{root} y se vea que falla, (c) se
ejecute con \texttt{sudo} y se vea que funciona, (d) se cambie el
dueño del propio programa a \texttt{root} y se vea qué pasa, (e) se
encienda el bit SUID y se ejecute como usuario regular --- explicando
en cada paso por qué.

\subsection*{Análisis: qué hace el bit SUID}

Cuando un usuario ejecuta un binario, el SO crea un proceso cuyo
\texttt{ruid} y \texttt{euid} son ambos iguales al UID del usuario.
Pero si el binario tiene el bit \textbf{SUID} encendido, el SO
asigna al proceso un \texttt{euid} \emph{igual al dueño del binario}.
En símbolos, al hacer \texttt{exec} de un programa con SUID:

\quad ruid\_nuevo = ruid\_actual\\
\quad euid\_nuevo = owner\_del\_programa

\begin{importante}
SUID es el mecanismo histórico de UNIX para permitir que un programa
sin privilegios delegue una acción privilegiada al kernel a través de
un binario que pertenece a un usuario privilegiado. El ejemplo
canónico es \texttt{/usr/bin/passwd}, que pertenece a \texttt{root}
con SUID encendido: lo ejecuta cualquier usuario pero corre como
\texttt{root}, lo cual le permite escribir en \texttt{/etc/shadow}.
\end{importante}

\subsection*{Programa extendido}

\begin{lstlisting}[style=cstyle]
#include <stdio.h>
#include <unistd.h>
#include <fcntl.h>
#include <sys/types.h>
#include <pwd.h>

int main(void) {
    uid_t ruid = getuid();
    uid_t euid = geteuid();
    struct passwd *pw_r = getpwuid(ruid);
    struct passwd *pw_e = getpwuid(euid);

    printf("ruid=%d (%s), euid=%d (%s)\n",
           ruid, pw_r ? pw_r->pw_name : "?",
           euid, pw_e ? pw_e->pw_name : "?");

    /* Abrir el archivo creado en el ejercicio 8 */
    int fd = open("test.txt", O_RDWR);
    if (fd < 0) {
        perror("open test.txt");
        return 1;
    }
    printf("Archivo abierto OK (fd=%d)\n", fd);
    close(fd);
    return 0;
}
\end{lstlisting}

Compilación:

\begin{lstlisting}[style=shellstyle]
$ gcc -Wall -o whoami_open whoami_open.c
\end{lstlisting}

Estado inicial: \texttt{test.txt} es del usuario regular con permisos
$600$, y \texttt{whoami\_open} también es del usuario regular con
permisos $755$.

\subsection*{10.a --- Abrir el archivo: OK}

\begin{lstlisting}[style=shellstyle]
$ ./whoami_open
ruid=1000 (tomi), euid=1000 (tomi)
Archivo abierto OK (fd=3)
\end{lstlisting}

\textbf{Por qué funciona.} El \texttt{euid} del proceso es $1000$
(\texttt{tomi}), el archivo es de \texttt{tomi} con permisos $600$
(\texttt{rw} para el dueño). El kernel chequea
``\texttt{euid == st\_uid \&\& (mode \& S\_IWUSR)}'' y deja pasar.

\subsection*{10.b --- Cambiar el dueño del archivo a \texttt{root}}

\begin{lstlisting}[style=shellstyle]
$ sudo chown root test.txt
$ ls -l test.txt
-rw-------  1 root staff  10 ... test.txt
\end{lstlisting}

\subsection*{10.c --- Repetir 10.a: ahora falla}

\begin{lstlisting}[style=shellstyle]
$ ./whoami_open
ruid=1000 (tomi), euid=1000 (tomi)
open test.txt: Permission denied
\end{lstlisting}

\textbf{Por qué falla.} El \texttt{euid} sigue siendo \texttt{tomi},
pero el archivo es de \texttt{root} con permisos $600$ (sólo el dueño
puede leer/escribir). El kernel rechaza con
\texttt{EACCES (Permission denied)}.

\subsection*{10.d --- Ejecutar con \texttt{sudo}: funciona}

\begin{lstlisting}[style=shellstyle]
$ sudo ./whoami_open
ruid=0 (root), euid=0 (root)
Archivo abierto OK (fd=3)
\end{lstlisting}

\textbf{Por qué funciona.} \texttt{sudo} ejecuta el comando con
\texttt{ruid = euid = 0} (\texttt{root}). El archivo es de
\texttt{root}, así que el chequeo del kernel pasa.

\begin{idea}
\texttt{sudo} es a su vez un binario \textbf{SUID-root}: pertenece a
root con bit SUID encendido. Por eso cualquier usuario lo ejecuta y se
le otorga (temporalmente, y previa autenticación) la identidad de
root. Lo que estamos viendo en 10.d es la consecuencia de un programa
SUID-root muy famoso.
\end{idea}

\subsection*{10.e --- Cambiar el dueño del programa a \texttt{root}}

\begin{lstlisting}[style=shellstyle]
$ sudo chown root whoami_open
$ ls -l whoami_open
-rwxr-xr-x  1 root staff  16400 ... whoami_open
\end{lstlisting}

El programa ahora es de root, pero conserva permisos $755$, así que el
usuario regular sigue pudiendo ejecutarlo.

\subsection*{10.f --- Ejecutar sin \texttt{sudo}: sigue fallando}

\begin{lstlisting}[style=shellstyle]
$ ./whoami_open
ruid=1000 (tomi), euid=1000 (tomi)
open test.txt: Permission denied
\end{lstlisting}

\textbf{Por qué falla.} Aunque el \emph{programa} sea de root, al
ejecutarlo el sistema le asigna al proceso un \texttt{euid} igual al
del usuario que lo invocó (\texttt{tomi}). Cambiar el dueño del
programa por sí solo \emph{no} cambia el \texttt{euid} del proceso.

\subsection*{10.g --- Con \texttt{sudo}: vuelve a funcionar}

\begin{lstlisting}[style=shellstyle]
$ sudo ./whoami_open
ruid=0 (root), euid=0 (root)
Archivo abierto OK (fd=3)
\end{lstlisting}

Misma razón que 10.d.

\subsection*{10.h --- Encender el bit SUID en el programa}

\begin{lstlisting}[style=shellstyle]
$ sudo chmod u+s whoami_open
$ ls -l whoami_open
-rwsr-xr-x  1 root staff  16400 ... whoami_open
\end{lstlisting}

\textbf{Lectura del nuevo permiso.} La letra \texttt{s} en lugar de
\texttt{x} en la posición del owner indica que el bit SUID está
encendido. (Si el archivo tuviera SUID pero no \texttt{x}, aparecería
una \texttt{S} mayúscula.)

En octal, SUID es el bit $4000$:

\quad $4755$ = SUID + \texttt{rwxr-xr-x}

\subsection*{10.i --- Ejecutar como usuario regular: ahora funciona}

\begin{lstlisting}[style=shellstyle]
$ ./whoami_open
ruid=1000 (tomi), euid=0 (root)        <- !!!
Archivo abierto OK (fd=3)
\end{lstlisting}

\textbf{Por qué funciona.} El kernel, al hacer \texttt{exec} del
binario SUID:

\begin{itemize}[leftmargin=*,nosep]
    \item Setea \texttt{ruid = 1000} (el que lanzó el proceso).
    \item \textbf{Setea \texttt{euid = 0}} (el dueño del programa).
\end{itemize}

A partir de ahí, cuando el programa llama \texttt{open("test.txt",
O\_RDWR)}, el kernel chequea con \texttt{euid = 0}: pasa el chequeo
porque \texttt{root} puede abrir cualquier cosa.

\subsection*{Resumen y diagrama del juego de UIDs}

\begin{center}
\small
\begin{tabular}{@{}llllll@{}}
\toprule
Paso & Dueño \texttt{test.txt} & Dueño \texttt{prog} & SUID? & ejec.\ con \texttt{sudo}? & resultado \\
\midrule
10.a & tomi & tomi & no & no  & OK \\
10.c & root & tomi & no & no  & \texttt{EACCES} \\
10.d & root & tomi & no & sí  & OK \\
10.f & root & root & no & no  & \texttt{EACCES} \\
10.g & root & root & no & sí  & OK \\
10.i & root & root & \textbf{sí} & no & \textbf{OK} \\
\bottomrule
\end{tabular}
\end{center}

\begin{center}
\begin{tikzpicture}[
    every node/.style={font=\ttfamily\small},
    box/.style={draw,rounded corners,minimum width=3.6cm,minimum height=0.7cm,align=center}
]
\node[box,fill=io1!20] (a) at (0,3) {usuario \texttt{tomi}};
\node[box,fill=io2!20] (b) at (0,1.7) {\texttt{exec("/usr/bin/prog")}};
\node[box,fill=io3!20] (c) at (0,0.4) {kernel chequea SUID};
\node[box,fill=ioshared!50] (d) at (0,-0.9) {ruid=tomi, euid=root};
\node[box,fill=iofile!50] (e) at (0,-2.2) {syscalls validadas con euid=root};

\draw[-Latex] (a) -- (b);
\draw[-Latex] (b) -- (c);
\draw[-Latex] (c) -- (d);
\draw[-Latex] (d) -- (e);
\end{tikzpicture}
\end{center}

\begin{importante}
\textbf{SUID es peligroso}. Un programa SUID-root con un bug que
permita ejecución arbitraria (buffer overflow, inyección de
comandos, race condition) le da \texttt{root} al atacante. Por eso
los binarios SUID-root del sistema son extremadamente auditados
(\texttt{passwd}, \texttt{sudo}, \texttt{su}, \texttt{ping} con
sockets raw, etc.) y se reducen al mínimo posible. Encontrarlos en un
sistema:

\begin{lstlisting}[style=shellstyle]
$ find / -perm -u+s -type f 2>/dev/null
\end{lstlisting}
\end{importante}

\textbf{Conclusión del ejercicio.} La secuencia ilustra los tres
elementos que el kernel mira al chequear permisos sobre un archivo:
(1) el \texttt{euid} del proceso (no el ruid), (2) el dueño y los
permisos del archivo (no del programa), y (3) si el binario tiene el
bit SUID, qué hace al \texttt{exec} (setear \texttt{euid} al dueño del
binario). Estos tres puntos en interacción son la base de la
seguridad clásica de UNIX.

% =====================================================================
\section{Conclusiones}
% =====================================================================

El práctico recorrió, ejercicio a ejercicio, las dos grandes capas del
subsistema de E/S que ofrece UNIX al programador:

\begin{itemize}[leftmargin=*]
    \item Los \textbf{Ejercicios 1 a 4} se concentraron en el
        \emph{lado bajo}, donde los dispositivos físicos aparecen como
        archivos especiales en \texttt{/dev}. El primero ejercitó la
        identificación de tipos de dispositivo en la salida de
        \texttt{ls}; el segundo y el tercero usaron pseudo-dispositivos
        (\texttt{/dev/zero} y \texttt{/dev/urandom}) como fuentes de
        datos; el cuarto exploró el flujo completo de construcción de
        un \emph{driver} como módulo cargable del kernel, con su
        comunicación con el RTC vía \texttt{ioctl} sobre
        \texttt{/dev/rtc0}.

    \item Los \textbf{Ejercicios 5 y 6} subieron a la abstracción del
        \emph{filesystem}: el 5 contrastó las dos representaciones
        canónicas (FAT centralizada vs.\ inodos descentralizados) y
        dejó en claro cuál es la diferencia conceptual (dónde vive la
        metadata y la cadena de bloques de cada archivo); el 6 mostró
        el patrón \texttt{readdir + stat} para inspeccionar entradas
        de un directorio desde C.

    \item Los \textbf{Ejercicios 7 a 10} cubrieron usuarios y control
        de acceso. El 7 fue el lado descriptivo (cómo el sistema
        guarda y expone usuarios y grupos en Linux y macOS); el 8, la
        mecánica concreta de los permisos UGO y octal sobre archivos
        y directorios, con énfasis en el bit \texttt{x} en
        directorios; el 9, la distinción entre \texttt{ruid} y
        \texttt{euid}; y el 10, el clímax del bloque: el bit
        \texttt{SUID}, que es la pieza que hace posible que un proceso
        no privilegiado obtenga, durante la ejecución de un binario
        cuidadosamente diseñado, los privilegios del dueño de ese
        binario.
\end{itemize}

\begin{idea}
La idea que une todo el práctico es que \textbf{en UNIX, todo es un
archivo}, y \textbf{para todo archivo hay un dueño y permisos}. La
combinación de esas dos abstracciones simples ---inodo + permisos
UGO--- es lo que permite que el mismo subsistema sirva indistintamente
para acceder a discos, dispositivos, sockets, pipes y memoria
compartida, con un control de acceso uniforme.
\end{idea}

\end{document}
